% Part: normal-modal-logic
% Chapter: axioms-systems
% Section: logics-proofs

\documentclass[../../../include/open-logic-section]{subfiles}

\begin{document}

\olfileid[hi]{nml}{prf}{prf}

\olsection{\usetoken{P}{derivation} और मोडल तंत्र}

पहले हम परिभाषित करते हैं कि सामान्य मोडल तर्कशास्त्रों के लिए
!!{derivation} क्या है। मोटे तौर पर, !!{derivation} !!{formula}ों का ऐसा
अनुक्रम है जिसमें प्रत्येक !!{element} या तो अनेक \emph{अभिगृहीतों} में से
किसी एक का (प्रतिस्थापन-दृष्टांत) है, या कुछ अनुमान-नियमों में से किसी नियम
से पूर्ववर्ती !!{element}ों से निकलता है। सामान्य मोडल तर्कशास्त्रों के लिए
सर्वसत्यताओं\iftag{prvDiamond}{, \Ax{K} और~\Dual{}}{ और~\Ax{K}} के सभी
दृष्टांत अभिगृहीत माने जाते हैं। इससे सबसे छोटा सामान्य मोडल तर्कशास्त्र,
मोडल तंत्र~$\Log{K}$, मिलता है। अन्य तंत्र पाने के लिए हम अतिरिक्त अभिगृहीत
जोड़ सकते हैं। नियम सदैव मोडस पोनेन्स~\MP{} और आवश्यकीकरण~\Nec{} हैं।

\begin{defn}
  कोई मोडल तंत्र $\Log{K} !A_1 \dots !A_n$ और !!{formula}~$!B$ दिए होने पर
  हम कहते हैं कि $!B$, $\Log{K} !A_1 \dots !A_n$ में
  \emph{!!{derivable}} है, जिसे $\Log{K} !A_1 \dots !A_n \Proves !B$
  लिखते हैं, तभी जब ऐसे !!{formula} $!C_1$, \dots, $!C_k$ हों कि $!C_k=!B$
  और प्रत्येक $!C_i$ या तो कोई सर्वसत्यतात्मक दृष्टांत हो, या $\Ax{K}$,
  \iftag{prvDiamond}{$\Dual$,}{} $!A_1$, \dots, $!A_n$ में से किसी का
  दृष्टांत हो, या नियम \MP{} अथवा~\Nec{} से पूर्ववर्ती !!{formula}ों से निकले।
\end{defn}

निम्न प्रतिज्ञप्ति हमें $!B$ की कोई $\Sigma$-!!{derivation} प्रस्तुत करके
$!B \in \Sigma$ दिखाने देती है।

\begin{prop}
  $\Log{K} !A_1 \dots !A_n = \Setabs{!B}{\Log{K} !A_1 \dots !A_n \Proves !B}$.
\end{prop}

\begin{proof}
  हम !!{derivation}यों की लंबाई पर आगमन से दिखाते हैं कि
  $\Setabs{!B}{\Log{K} !A_1 \dots !A_n \Proves !B} \subseteq
  \Log{K} !A_1 \dots !A_n$.

  यदि $!B$ की !!{derivation} की लंबाई~$1$ है, तो उसमें एक ही !!{formula} है।
  वह !!{formula} \MP{} या \Nec{} से पूर्ववर्ती सूत्रों से नहीं निकल सकता,
  इसलिए वह कोई सर्वसत्यतात्मक दृष्टांत, \Ax{K},
  \iftag{prvDiamond}{$\Dual$,}{} या $!A_1$, \dots,~$!A_n$ में से किसी का
  दृष्टांत होना चाहिए। $\Log{K}!A_1\dots!A_n$ में ये भी हैं, अतः
  $!B \in \Log{K}!A_1 \dots !A_n$।

  यदि $!B$ की !!{derivation} की लंबाई~$>1$ है, तो $!B$ इसके अतिरिक्त
  !!{derivation} की अंतिम पंक्ति से पहले आने वाले !!{formula}ों से \MP{} या
  \Nec{} द्वारा मिल सकता है। यदि $!B$, $!C$ और $!C \lif !B$ से (\MP{} द्वारा)
  निकलता है, तो आगमन-परिकल्पना से $!C$ और
  $!C \lif !B \in \Log{K}!A_1 \dots !A_n$। प्रत्येक मोडल तर्कशास्त्र मोडस
  पोनेन्स के अधीन संवृत है, अतः $!B \in \Log{K}!A_1 \dots !A_n$। यदि
  $!B \equiv \Box !C$, $!C$ से \Nec{} द्वारा निकलता है, तो आगमन-परिकल्पना
  से $!C \in \Log{K}!A_1 \dots !A_n$। प्रत्येक सामान्य मोडल तर्कशास्त्र
  $\Nec$ के अधीन संवृत है, अतः $!B \in \Log{K}!A_1\dots!A_n$।

  विलोम समावेशन यह दिखाने से मिलता है कि
  $\Sigma = \Setabs{!B}{\Log{K} !A_1 \dots !A_n \Proves !B }$
  ऐसा सामान्य मोडल तर्कशास्त्र है जिसमें $!A_1$, \dots, $!A_n$ के सभी
  दृष्टांत हैं, और यह देखने से कि परिभाषा से $\Log{K} !A_1 \dots !A_n$
  ऐसा सबसे छोटा तर्कशास्त्र है।
  \begin{enumerate}
    \item प्रत्येक सर्वसत्यता~$!B$ सर्वसत्यतात्मक दृष्टांत है, अतः
      $\Log{K}!A_1\dots!A_n \Proves !B$; फलतः
      $\Sigma$ में सभी सर्वसत्यताएँ हैं।
    \item यदि $\Log{K}!A_1\dots!A_n \Proves !C$ और
      $\Log{K}!A_1\dots!A_n \Proves !C \lif !B$, तो
      $\Log{K}!A_1\dots!A_n \Proves !B$: $!C$ की !!{derivation} को
      $!C \lif !B$ की !!{derivation} से मिलाएँ और पंक्ति~$!B$ जोड़ें। अंतिम
      पंक्ति \MP{} से उचित है। अतः $\Sigma$ मोडस पोनेन्स के अधीन संवृत है।
    \item यदि $!B$ की कोई !!{derivation} है, तो $!B$ के प्रत्येक
      प्रतिस्थापन-दृष्टांत की भी !!{derivation} है: प्रतिस्थापन को
      !!{derivation} के प्रत्येक !!{formula} पर लागू करें। (अभ्यास:
      !!{derivation} की लंबाई पर आगमन से सिद्ध करें कि परिणाम भी सही
      !!{derivation} है।) अतः $\Sigma$ एकरूप प्रतिस्थापन के अधीन संवृत है।
      (अब हमने स्थापित किया है कि $\Sigma$ मोडल तर्कशास्त्र की सभी शर्तें
      संतुष्ट करता है।)
    \item हमारे पास $\Log{K}!A_1\dots!A_n \Proves \Ax{K}$, अतः
      $K \in \Sigma$।
    \tagitem{prvDiamond}{हमारे पास
      $\Log{K}!A_1\dots!A_n \Proves \Dual$, अतः $\Dual \in \Sigma$।}{}
    \item यदि $\Log{K}!A_1\dots!A_n \Proves !C$, तो अतिरिक्त
      पंक्ति~$\Box !C$, \Nec{} से उचित है। फलतः $\Sigma$, \Nec{} के अधीन
      संवृत है। अतः $\Sigma$ सामान्य है।
  \end{enumerate}
\end{proof}

\end{document}
